\section{Roots of Nonlinear Equations}

\subsection{Some observations about the convergence rate}\label{sec:convergence-rate}

\paragraph{Convergence rate:}
By definition, an iterative method has an order of convergence $q$ and an
asymptotic error constant $C$ if the sequence of errors $\epsilon_k = x_k - r$
satisfies:
\begin{equation*}
	|\epsilon_{k+1}| \approx C |\epsilon_k|^q
\end{equation*}
If we don't know the exact root $r$, we can define the step difference $d_k$ as
the distance between two consecutive approximations:
\begin{equation*}
	d_k = |x_{k+1} - x_k|
\end{equation*}
We can rewrite this by subtracting and adding the true root $r$:
\begin{equation*}
	d_k = |(x_{k+1} - r) - (x_k - r)| = |\epsilon_{k+1} - \epsilon_k|
\end{equation*}
Because the algorithm is converging, the error at the next step is much smaller
than the error at the current step (i.e., $|\epsilon_{k+1}| \ll |\epsilon_k|$).
Therefore, $\epsilon_{k+1}$ becomes negligible compared to $\epsilon_k$.
We can safely approximate:
\begin{equation*}
	d_k \approx |\epsilon_k|
\end{equation*}
Since $d_k \approx |\epsilon_k|$ and $d_{k+1} \approx |\epsilon_{k+1}|$, we can
substitute these into our original theoretical definition:
\begin{equation*}
	d_{k+1} \approx C (d_k)^q
\end{equation*}
Now, take the natural logarithm of both sides to bring the exponent
$q$ down:
\begin{equation*}
	\ln(d_{k+1}) \approx \ln(C) + q \ln(d_k)
\end{equation*}
This equation has two unknowns: $C$ and $q$. To eliminate the constant $C$, we
write the exact same equation for the previous step $k-1$:
\begin{equation*}
	\ln(d_k) \approx \ln(C) + q \ln(d_{k-1})
\end{equation*}
Subtract the second equation from the first equation. This completely
eliminates $\ln(C)$:
\begin{equation*}
	\ln(d_{k+1}) - \ln(d_k) \approx q \ln(d_k) - q \ln(d_{k-1})
\end{equation*}
Factor out $q$ on the right side:
\begin{equation*}
	\ln(d_{k+1}) - \ln(d_k) \approx q[ \ln(d_k) - \ln(d_{k-1}) ]
\end{equation*}
Using the logarithm quotient rule ($\ln(a) - \ln(b) = \ln(a/b)$), we rewrite
this as:
\begin{equation*}
	\ln\left(\frac{d_{k+1}}{d_k}\right) \approx q \ln\left(\frac{d_k}{d_{k-1}}\right)
\end{equation*}
Finally, divide to isolate $q$:
\begin{equation*}
	q \approx \frac{\ln(d_{k+1} / d_k)}{\ln(d_k / d_{k-1})}
\end{equation*}
This equation allows us to calculate the convergence rate without knowing the
exact root of an equation.

\paragraph{Asymptotic error constant:}
The strict mathematical definition of the asymptotic error constant $C$ is
defined by the limit:
\begin{equation*}
	\lim_{k \to \infty} \frac{|\epsilon_{k+1}|}{|\epsilon_k|^q} = C
\end{equation*}
Here, the error is $\epsilon_k = x_k - r$.
For a sufficiently large $k$ (meaning we have run the algorithm for a few steps
and are close to the root), we can drop the limit and write this as an
approximation:
\begin{equation*}
	C \approx \frac{|\epsilon_{k+1}|}{|\epsilon_k|^q} = \frac{|x_{k+1} - r|}{|x_k - r|^q}
\end{equation*}
To compute this practically, we run into the same problem: we don't know $r$.
However, iterative algorithms are typically run until the difference between
steps drops below machine precision.

Therefore, the final approximation the algorithm produces, $x_{final}$, is
virtually indistinguishable from the exact root $r$ for our precision levels.
We make the substitution $r \approx x_{final}$:
\begin{equation*}
	C \approx \frac{|x_{k+1} - x_{final}|}{|x_k - x_{final}|^q}
\end{equation*}

\subsection{Bisection Method}

\begin{exercise}
	\begin{enumerate}
		\item Use the bisection method to find the solutions for $x^3 - 7x^2 +
			      14x - 6 = 0$ on $[0,\, 1]$.
		\item Study the convergence of the bisection method to the root found in (1).
		      You can take as an approximation for the root $r = m_{n_\text{max}}$
		      where $n_\text{max}$ is the index corresponding to the last bisection
		      iteration. What is its order of convergence $q$? Can you estimate the
		      asymptotic error constant $C$? [Hint: Study $a_{n} = -\log_{10}\lvert
				      x_{n} - r \rvert$ as a function of $n$.]
	\end{enumerate}
\end{exercise}

\begin{solution}
	The iterative process, to solve the root-finding problem for the non-linear
	function $f(x) = x^3 - 7x^2 + 14x - 6$, was designed to terminate when the
	interval width, representing the maximum absolute error, dropped below a
	strictly defined tolerance of $\texttt{tol} = 10^{-14}$.

	To empirically determine the order of convergence $q$ and the asymptotic
	error constant $C$, we define the absolute error at the $n$-th iteration as
	$\epsilon_n = |x_n - r|$, where $r$ is our terminal approximation. Assuming
	the error propagates according to the asymptotic relation $\epsilon_{n+1} \approx
		C \epsilon_n^q$, we can linearize the expression via a base-10 logarithm:
	\[
		\log_{10}|\epsilon_{n+1}| \approx q \log_{10}|\epsilon_n| + \log_{10}C
	\]
	A linear least-squares regression was applied to the sequence of coordinate
	pairs $(\log_{10}|\epsilon_n|, \log_{10}|\epsilon_{n+1}|)$. Because initial
	iterations often fall outside the asymptotic regime, we performed a
	sensitivity analysis by systematically truncating the first $k$ iterations
	(where $k \in \{1, 3, 5, \dots, 19\}$) to isolate the asymptotic behavior of
	the sequence's tail.

	\paragraph{Results}
	The bisection algorithm successfully converged to the root $r \approx
		0.5857864376269042$ after 47 iterations. The geometric progression of the
	root extraction is illustrated in Figure \ref{fig:bisection_exact}.

	\begin{center}
		\centering
		% Ensure path matches your standard inclusion macros
		\includegraphics[width=0.8\linewidth]{figures/4-2/Exact solution: x = 0.5857864376269042 with 47 iterations.pdf}
		\captionof{figure}{Function plot and the convergence of the bisection
			method to the exact root after 47 iterations.}
		\label{fig:bisection_exact}
	\end{center}

	The empirical evaluation of the order of convergence $q$ and error constant
	$C$ yielded varied results depending on the truncation of the initial
	estimates. A selection of these results is tabulated below:

	\begin{center}
		\centering
		\renewcommand{\arraystretch}{1.2}
		\begin{tabular}{cccc}
			\toprule
			Skipped Iterations ($k$) & Analyzed Tail & Convergence Order ($q$) & Error Constant ($C$) \\
			\midrule
			1                                 & 46                     & 0.9789                           & 0.3431                        \\
			5                                 & 42                     & 0.9744                           & 0.3078                        \\
			9                                 & 38                     & 0.9688                           & 0.2616                        \\
			15                                & 32                     & 0.9558                           & 0.1822                        \\
			19                                & 28                     & 0.9391                           & 0.1165                        \\
			\bottomrule
		\end{tabular}
		\captionof{table}{Empirical convergence parameters $q$ and $C$ as a
			function of skipped initial iterations.}
		\label{tab:convergence_params}
	\end{center}

	The log-log regression for the case where 1 initial iteration was skipped is
	visually represented in Figure \ref{fig:conv_analysis_1}, yielding $q
		\approx 0.979$ and $C \approx 0.343$.

	\begin{center}
		\centering
		\includegraphics[width=0.8\linewidth]{figures/4-2/conv-analysis-with-1-skipped.pdf}
		\captionof{figure}{Linear least-squares fit for convergence analysis,
			skipping the first initial iteration.}
		\label{fig:conv_analysis_1}
	\end{center}

	\paragraph{Discussion}
	The theoretical expectation for the bisection method is strict linear
	convergence, characterized by an order of $q = 1$ and an asymptotic error
	constant of $C = 0.5$, reflecting the halving of the search interval at each
	step. Our empirical findings ($q \approx 0.94 - 0.98$ and $C \approx 0.12 -
		0.34$) loosely approximate, but slightly underestimate, these theoretical
	markers.

	This discrepancy arises from the nature of the bisection method: while the
	maximum bound of the error decreases monotonically by exactly $1/2$
	at each step, the actual error $\epsilon_n = |x_n - r|$ does not
	strictly decrease monotonically. Depending on the geometry of the interval, a
	given approximation $x_n$ can occasionally land extremely close to the true
	root early on, introducing high variance and non-smooth progression in the
	$\log_{10}|\epsilon_n|$ sequence. Consequently, performing a least-squares fit
	on consecutive actual errors yields a noisy representation of the underlying
	linear upper bound envelope.

	\paragraph{Global Convergence Analysis}
	To further validate the linear convergence of the bisection method, we
	investigated the relationship between the iteration index $n$ and the
	logarithmic error measure $a_n = -\log_{10}|\epsilon_n|$, as suggested.
	Under the assumption of linear convergence, the error evolves as $\epsilon_n
		\approx \epsilon_0 C^n$. Taking the negative base-10 logarithm yields:
	\[
		a_n \approx -n \log_{10}(C) - \log_{10}(\epsilon_0)
	\]
	This implies that a plot of $a_n$ versus $n$ should yield a linear
	trajectory with a slope $S = -\log_{10}(C)$. This global perspective is
	numerically more robust than step-by-step ratios, as it averages the
	non-monotonic fluctuations inherent in the bisection method over the entire
	iterative history.

	\paragraph{Results \& Data Visualization}
	The bisection method converged to $r \approx 0.5857864376269042$ in 47
	iterations (Fig. \ref{fig:bisection_exact}). The global convergence analysis
	(Fig. \ref{fig:global_conv}) shows a highly linear trend ($R^2 \approx 1$).
	The calculated slope $S \approx 0.299$ corresponds to an empirical error
	constant of:
	\[
		C = 10^{-S} \approx 10^{-0.299} \approx 0.5000
	\]
	This result aligns perfectly with the theoretical expectation for the
	bisection method, where the uncertainty interval is halved ($C = 0.5$) at
	each step.

	\begin{center}
		\centering
		\includegraphics[width=0.8\linewidth]{figures/4-2/global-convergence-an.pdf}
		\captionof{figure}{Global convergence analysis: $a_n =
				-\log_{10}|\epsilon_n|$ as a function of the iteration index $n$. The
			linear fit confirms a steady increase in significant digits.}
		\label{fig:global_conv}
	\end{center}

	\paragraph{Discussion}
	The discrepancy noted in the previous step-by-step analysis, where $q$ and
	$C$ fluctuated based on the number of skipped iterations, is resolved by the
	global $a_n$ vs. $n$ approach. While the distance between the midpoint $m_n$
	and the true root $r$ can fluctuate non-monotonically (causing "noise" in
	the $\epsilon_{n+1}/\epsilon_n$ ratio), the maximum possible error is
	strictly halved.

	\paragraph{Conclusion}
	The bisection implementation proved highly stable, converging to a tight
	tolerance within the operational limits of double-precision arithmetic.
	While empirical analysis confirms the fundamentally linear nature of the
	convergence ($q \to 1$), it also highlights the difference between tracking
	strict theoretical error bounds versus highly variable actual sequence
	errors in deterministic interval-halving algorithms.

\end{solution}

\subsection{Newton Method}

\begin{exercise}
	For each of the functions below, do the following steps.

	\begin{itemize}
		\item Rewrite the equation into the standard form for rootfinding, $f(x)
			      = 0$, and compute $f'(x)$.
		\item Make a plot of $f$ over the given interval and determine how many
		      roots lie in the interval.
		\item Use Newton's method to find each root.
		\item Study the error in the Newton sequence to determine numerically
		      whether the convergence is roughly quadratic for the given root. [Hint:
				      Study $a_{n+1}/a_{n}$ with $a_n=-\log_{10}|x_n-r|$ to check the rate of
				      convergence $q$. Once you established it, try to estimate the asymptotic
				      error constant $C$.]

		      \begin{itemize}
			      \item $x^2=e^{-x}$, over $[-2,2]$

			      \item$2x = \tan x$, over $[-0.2,1.4]$

			      \item $e^{x+1}=2+x$, over $[-2,2]$
		      \end{itemize}
	\end{itemize}
\end{exercise}

\begin{solution}
	To apply Newton's method, we first rewrite the given equations into the
	standard root-finding form $f(x) = 0$ and compute their respective
	derivatives, $f'(x)$.

	\begin{enumerate}
		\item For $x^2 = e^{-x}$:
		      \begin{equation*}
			      f_1(x) = x^2 - e^{-x}, \quad f_1'(x) = 2x + e^{-x}.
		      \end{equation*}
		\item For $2x = \tan x$:
		      \begin{equation*}
			      f_2(x)  = 2x - \tan x, \quad f_2'(x) = 2 - \sec^2 x.
		      \end{equation*}
		\item For $e^{x+1} = 2+x$:
		      \begin{equation*}
			      f_3(x) = e^{x+1} - x - 2, \quad f_3'(x) = e^{x+1} - 1.
		      \end{equation*}
	\end{enumerate}

	\paragraph{Numerical Methodology}
	The core algorithm computes the Newton sequence via the recurrence relation $x_{n+1} = x_n - f(x_n)/f'(x_n)$.

	To evaluate the convergence rate $q$ and the asymptotic error constant $C$, we
	used the algorithm explained in the Section~\ref{sec:convergence-rate}.

	\paragraph{Results}
	Graphical analysis over the designated intervals (Figure
	\ref{fig:newton_roots}) confirms the existence of a single root within the
	domain for each function. The quantitative results of the Newton iterations
	are summarized in Table \ref{tab:newton_results}.

	\begin{center}
		\centering
		% Ensure the path matches your LaTeX document structure
		\includegraphics[width=\textwidth]{figures/4-3/plot-solution-newton-method.pdf}
		\captionof{figure}{Behavior of the objective functions $f_1(x)$, $f_2(x)$, and $f_3(x)$ over their respective intervals. The dashed line represents $y=0$, and the calculated roots are denoted by scatter points.}
		\label{fig:newton_roots}
	\end{center}

	\begin{center}
		\centering
		\captionof{table}{Numerical roots $r$, observed convergence rates $q$, and asymptotic error constants $C$ obtained via Newton's method in \texttt{Float64} precision.}
		\label{tab:newton_results}
		\renewcommand{\arraystretch}{1.2}
		\begin{tabular}{l c c c}
			\toprule
			Function          & Root Estimate ($r$) & Convergence Rate ($q$) & Error Constant ($C$) \\
			\midrule
			$f_1(x) = x^2 - e^{-x}$    & $0.703467422498$             & $2.0007$                        & $0.3978$                      \\
			$f_2(x) = 2x - \tan x$     & $1.165561185207$             & $2.0003$                        & $3.3956$                      \\
			$f_3(x) = e^{x+1} - x - 2$ & $-1.00000000372$             & $1.0113$                        & $0.7841$                      \\
			\bottomrule
		\end{tabular}
	\end{center}

	\paragraph{Discussion}
	The numerical results strongly align with theoretical expectations for Newton's method:
	\begin{itemize}
		\item Functions 1 and 2: Both exhibit a convergence rate of $q
			      \approx 2$. This numerically confirms quadratic convergence, which is
		      expected since the roots are simple (i.e., $f'(r) \neq 0$). The process
		      successfully converges to machine epsilon limits without numerical
		      instability.
		\item Function 3: The convergence rate degenerates to $q
			      \approx 1$ (linear convergence). This is analytically justified by
		      evaluating the exact root at $x = -1$. Observe that $f_3(-1) = 0$ and
		      $f'_3(-1) = 0$. Because the first derivative vanishes at the root, the
		      root has a multiplicity of $m=2$. Newton's method theoretically reverts
		      to linear convergence for multiple roots.
	\end{itemize}
	Furthermore, for $f_3(x)$, the evaluation of $x_{n+1}$ near the root
	requires dividing a very small number by another very small number
	($f(x_n)/f'(x_n) \to 0/0$). This induces catastrophic cancellation in
	standard \texttt{Float64} arithmetic. Consequently, our estimate of the root
	carries an absolute error on the order of $10^{-9}$ instead of reaching the
	standard $10^{-16}$ machine epsilon bound, and the asymptotic constant $C
		\approx 0.78$ deviates from the theoretical value of $C = 1 - 1/m = 0.5$.

	\paragraph{Conclusion}
	Newton's method provides robust and highly efficient (quadratic)
	root-finding capabilities for simple roots. However, the presence of
	multiple roots severely impacts both the asymptotic convergence rate and the
	ultimate numerical stability of the algorithm.
\end{solution}

\begin{exercise}
	Plot the function $f(x)=x^{-2} - \sin x$ on the interval $x \in [0.5,10]$.
	For each initial value $x_1=1,\, x_1=2,\,\ldots,\, x_1=7$, apply Newton's
	method to $f$, and make a table showing $x_1$ and the resulting root found
	by the method. In which case does the iteration converge to a root other
	than the one closest to it? Use the plot to explain why that happened.
	(Optional) Repeat the exercise including bracketing.
\end{exercise}

\begin{solution}
	To apply Newton's method, we require the first derivative of the objective
	function $f(x) = x^{-2} - \sin(x)$. Applying standard differentiation
	yields:
	\[
		f'(x) = -2x^{-3} - \cos(x).
	\]

	\paragraph{Results}
	The target function $f(x)$ exhibits oscillatory behavior bounded by a
	$1/x^2$ decay, leading to multiple roots within the interval $[0.5, 10]$.
	Figure \ref{fig:multiple_roots} displays the objective function and the
	roots successfully identified by the solver. Table
	\ref{tab:newton_multiple_roots} details the root converged upon for each
	initial guess $x_1$, alongside convergence metrics.

	\begin{center}
		\centering
		\includegraphics[width=\textwidth]{figures/4-3/multiple-newton-method-roots.pdf}
		\captionof{figure}{Plot of $f(x) = x^{-2} - \sin(x)$ over the interval $x \in [0.5, 10]$. The dashed line denotes $y=0$, and the scatter points represent the numerical roots found via Newton's method.}
		\label{fig:multiple_roots}
	\end{center}

	\begin{center}
		\centering
		\captionof{table}{Newton's method results for varying initial guesses.}
		\label{tab:newton_multiple_roots}
		\renewcommand{\arraystretch}{1.2}
		\begin{tabular}{c c c c c}
			\toprule
			Initial Guess ($x_1$) & Converged Root ($r$) & Closest Root & Conv. Rate ($q$) \\
			\midrule
			$1.0$                          & $1.06822$                     & $1.06822$             & $2.0024$                  \\
			$2.0$                          & $6.30832$                     & $3.03265$             & $2.0283$                  \\
			$3.0$                          & $3.03265$                     & $3.03265$             & $2.0221$                  \\
			$4.0$                          & $3.03265$                     & $3.03265$             & $1.9756$                  \\
			$5.0$                          & $9.41349$                     & $6.30832$             & $1.7087$                  \\
			$6.0$                          & $6.30832$                     & $6.30832$             & $2.0188$                  \\
			$7.0$                          & $6.30832$                     & $6.30832$             & $2.1579$                  \\
			\bottomrule
		\end{tabular}
	\end{center}

	\paragraph{Discussion}
	As observed in Table \ref{tab:newton_multiple_roots}, the iterations
	initialized at $x_1 = 2.0$ and $x_1 = 5.0$ converge to roots far away from
	their starting positions, completely missing their geometrically closest
	roots.

	This phenomenon is fundamentally tied to the geometry of the function at the
	initial guess, as seen in Figure \ref{fig:multiple_roots}. The Newton step
	size is inversely proportional to the derivative $f'(x_n)$.
	\begin{itemize}
		\item Case $x_1 = 2.0$: The function has a local minimum
		      slightly before $x=2$. At $x=2$, the derivative $f'(2)$ is relatively
		      small, making the tangent line shallow. Because $f(2)$ is negative and
		      the slope is gently positive, the tangent line projects far across the
		      $x$-axis to the right, stepping over the nearby root at $r \approx 3.03$
		      and landing directly into the basin of attraction for the root at $r
			      \approx 6.31$.
		\item Case $x_1 = 5.0$: Similarly, $x=5$ sits near a local
		      maximum (which occurs near $3\pi/2 \approx 4.71$). The shallow negative
		      slope at $x=5$ projects the next guess far to the right, skipping over
		      the closest root $r \approx 6.31$ and instead converging to $r \approx
			      9.41$.
	\end{itemize}

	\paragraph{Conclusion}
	These results highlight the primary vulnerability of Newton's method: its
	lack of global convergence. For highly oscillatory functions like $x^{-2} -
		\sin(x)$, local extrema create "blind spots" where vanishing derivatives
	($f'(x) \to 0$) hurl the numerical sequence into entirely different basins
	of attraction.
\end{solution}

\begin{exercise}
	The most easily observed properties of the orbit of a celestial body around
	the Sun are the period $\tau$ and the elliptical eccentricity $\epsilon$.
	(A circle has $\epsilon=0$.) From these, it is possible to find at any time
	$t$ the true anomaly $\theta(t)$. This is the angle between the direction
	of periapsis (the point in the orbit closest to the Sun) and the current
	position of the body, as seen from the main focus of the ellipse (the
	location of the Sun). This is done through

	\begin{equation}
		\tan \frac{\theta}{2} = \sqrt{\frac{1+\epsilon}{1-\epsilon}}\,
		\tan \frac{\psi}{2}\,,
		\label{eq:eq-kepler1}
	\end{equation}

	where the eccentric anomaly $\psi(t)$ satisfies Kepler's equation:

	\begin{equation}
		\psi - \epsilon \sin \psi - \frac{2\pi t}{\tau} = 0\,.
		\label{eq:eq-kepler2}
	\end{equation}

	\ref{eq:eq-kepler2} must be solved numerically to find $\psi(t)$, and
	then \ref{eq:eq-kepler1} can be solved analytically to find $\theta(t)$.

	The asteroid Eros has $\tau=1.7610$ years and $\epsilon=0.2230$. Using
	Newton's method for \ref{eq:eq-kepler2}, make a plot of $\theta(t)$ for 100
	values of $t$ between $0$ and $\tau$, which is one full orbit.
\end{exercise}

\begin{solution}
	To compute the true anomaly $\theta(t)$ over time, we must first isolate the
	eccentric anomaly $\psi(t)$ by finding the roots of Kepler's equation at
	discrete time steps $t$. We define our objective function and its first
	derivative with respect to $\psi$ as follows:
	\[
		\begin{aligned}
			f(\psi)  & = \psi - \epsilon \sin \psi - \frac{2\pi t}{\tau}, \\
			f'(\psi) & = 1 - \epsilon \cos \psi.
		\end{aligned}
	\]
	Because the physical eccentricity of an ellipse satisfies $0 \le \epsilon <
		1$, we observe that $f'(\psi) \ge 1 - \epsilon > 0$ for all $\psi$. This
	implies that $f(\psi)$ is strictly monotonically increasing and contains
	exactly one unique real root for any given time $t$.

	\paragraph{Numerical Methodology}
	The orbital parameters for the asteroid Eros were defined as $\tau = 1.7610$
	years and $\epsilon = 0.2230$. The time domain $t \in [0, \tau]$ was
	discretized into 100 evenly spaced intervals.

	For each time step $t_i$, Newton's method was utilized to converge upon
	$\psi(t_i)$. Following convergence, the true anomaly $\theta(t_i)$ was
	computed via a rearrangement of Equation~\ref{eq:eq-kepler1}:
	\[
		\theta(t) = 2 \arctan\left( \sqrt{\frac{1+\epsilon}{1-\epsilon}} \tan\left(\frac{\psi(t)}{2}\right) \right).
	\]

	\paragraph{Results}
	The Newton iterations converged successfully for all 100 time steps. Figure
	\ref{fig:true_anomaly} illustrates the resulting true anomaly $\theta(t)$
	over the course of exactly one orbital period $\tau$.

	\begin{center}
		\centering
		\includegraphics[width=0.85\textwidth]{figures/4-3/full-orbit.pdf}
		\captionof{figure}{Evolution of the true anomaly $\theta(t)$ over one
			full orbital period ($\tau=1.7610$ years) for the asteroid Eros
			($\epsilon = 0.2230$).}
		\label{fig:true_anomaly}
	\end{center}

	\paragraph{Discussion}
	The plot of $\theta(t)$ demonstrates a monotonic increase from $0$ to $2\pi$
	radians, satisfying the boundary conditions of a closed orbit. Notably, the
	curve is not perfectly linear. This non-linearity is a direct macroscopic
	manifestation of Kepler's Second Law of Planetary Motion (equal areas in
	equal times). Because $\epsilon = 0.2230 > 0$, the asteroid travels
	significantly faster at periapsis ($\theta=0, 2\pi$) and slower at apoapsis
	($\theta=\pi$).

	\paragraph{Conclusion}
	Newton's method proves to be a robust algorithm for solving Kepler's equation.
\end{solution}

\subsection{Interpolations methods}

\begin{exercise}\label{ex:4.3.2}
	For each of the functions below, do the following steps.
	\begin{itemize}
		\item Rewrite the equation into the standard form for rootfinding, $f(x)
			      = 0$.
		\item Make a plot of $f$ over the given interval and determine how many
		      roots lie in the interval.
		\item For each root determine an interval that encloses the given root.
		      Then use the secant method, starting with the endpoints of the
		      enclosing interval, to find each root.
		\item For each root, use the errors in the secant sequence to determine
		      numerically whether the convergence is apparently between linear and
		      quadratic.
	\end{itemize}
	\begin{align*}
		x^2     & = e^{-x}, \quad \text{over } [-2,2]     \\
		2x      & = \tan x, \quad \text{over } [-0.2,1.4] \\
		e^{x+1} & = 2+x, \quad \text{over } [-2,2]
	\end{align*}

\end{exercise}

\begin{solution}
	We begin by expressing each equation in the standard root-finding form $f(x) = 0$:
	\begin{align*}
		f_1(x) & = x^2 - e^{-x} = 0, \quad    & x \in [-2, 2]    \\
		f_2(x) & = 2x - \tan(x) = 0, \quad    & x \in [0.2, 1.4] \\
		f_3(x) & = e^{x+1} - 2 - x = 0, \quad & x \in [-2, 2]
	\end{align*}
	Visual inspection of the functions (Figure~\ref{fig:secant_plots}) reveals
	that each interval contains exactly one valid root. The boundaries of the given intervals were
	utilized as the initial guesses, $x_0$ and $x_1$.

	\begin{center}
		\centering
		\includegraphics[width=\textwidth]{figures/4-4/secant-method-plot.pdf}
		\captionof{figure}{Plots of the functions $f_1(x)$, $f_2(x)$, and $f_3(x)$ over
			their respective intervals. The dashed line represents $y=0$, and the
			calculated roots are marked with points.}
		\label{fig:secant_plots}
	\end{center}

	\paragraph{Numerical Results}
	The secant method successfully converged for all three functions. To
	evaluate the order of convergence $q$, two distinct algorithms were
	employed. Algorithm 1 utilises a linear least-squares fit on the
	logarithmic errors, whereas Algorithm 2 computes the local
	convergence rate using successive differences (see
	Section~\ref{sec:convergence-rate}). The computed roots $x^*$ and the
	estimated convergence rates are summarised in
	Table~\ref{tab:secant_results}.

	\begin{center}
		\centering
		\captionof{table}{Roots found via the secant method and corresponding
			estimated orders of convergence ($q$). The theoretical expected rate for simple
			roots is $q \approx 1.618$.}
		\label{tab:secant_results}
		\renewcommand{\arraystretch}{1.2}
		\begin{tabular}{l c S[table-format=2.6] S[table-format=2.3] c}
			\toprule
			Function & Interval & {Root ($x^*$)} & {$q$ (Alg 1)} & {$q$ (Alg 2)} \\
			\midrule
			$f_1(x)$          & $[-2, 2]$         & 0.703467                & 1.606                  & 2.368                  \\
			$f_2(x)$          & $[0.2, 1.4]$      & 1.165561                & 0.073                  & 39.283                 \\
			$f_3(x)$          & $[-2, 2]$         & -1.000000               & 1.016                  & {\texttt{NaN}}         \\
			\bottomrule
		\end{tabular}
	\end{center}

	\paragraph{Discussion on Convergence Rates}
	The standard secant method yields superlinear convergence with $q = \frac{1
			+ \sqrt{5}}{2} \approx 1.618$ for simple roots.
	\begin{itemize}
		\item For $f_1(x)$, Algorithm 1 accurately captures this superlinear
		      behaviour ($q \approx 1.606$), closely aligning with theoretical
		      expectations.
		\item For $f_2(x)$, both algorithms fail to produce the expected $q$.
		      This is likely due to the extremely rapid convergence (requiring very
		      few iterations) combined with the steep gradient of $\tan(x)$, which
		      prevents the generation of a sufficiently long asymptotic error sequence
		      before machine precision limits are reached.
		\item For $f_3(x)$, the algorithm drops to linear convergence ($q
			      \approx 1.016$). This perfectly matches numerical theory: evaluating the
		      derivative $f_3'(x) = e^{x+1} - 1$ at $x = -1$ yields $0$. Therefore, $x
			      = -1$ is a multiple root (multiplicity $m=2$), and root-finding
		      methods strictly degrade to linear convergence in such cases.
	\end{itemize}

	\paragraph{Algorithmic Robustness: Algorithm 1 vs. Algorithm 2}
	As observed in Table \ref{tab:secant_results}, Algorithm 1 provides
	significantly more reliable estimates of $q$ compared to Algorithm 2. This
	discrepancy arises from how floating-point arithmetic interacts with the
	underlying mathematics of each algorithm.

	Algorithm 2 attempts to estimate $q$ locally using the ratio of successive differences:
	\[
		q_k \approx \frac{\ln(d_{k+1} / d_k)}{\ln(d_k / d_{k-1})}, \quad \text{where } d_k = |x_k - x_{k-1}|
	\]
	As the sequence converges, $d_k$ approaches the machine epsilon.
	Subtracting two nearly identical floating-point numbers triggers
	catastrophic cancellation, severely amplifying numerical noise.
	Despite implementing a truncation threshold ($d_i < 10^{-7}$), the sequence
	becomes too short or highly erratic, leading to wild overestimations (e.g.,
	$q \approx 39.2$ for $f_2$) or total failure (e.g., \texttt{NaN} for $f_3$
	due to early termination or division by zero).

	Conversely, Algorithm 1 applies a least-squares linear regression to the
	global error history:
	\[
		\log_{10} \epsilon_{n+1} \approx q \log_{10} \epsilon_n + \log_{10} C
	\]
	By formulating the problem as an overdetermined system and solving it via
	matrices, Algorithm 1 utilizes multiple data points. This least-squares
	approach inherently acts as a statistical filter, smoothing out the
	high-frequency quantization noise caused by floating-point arithmetic.
	Consequently, Algorithm 1 remains robust and stable even when local errors
	approach the threshold of numerical precision.

	\paragraph{Conclusion}
	Overall, the secant method successfully located the roots of all three
	nonlinear equations, demonstrating its effectiveness as a derivative-free
	root-finding technique. The numerical experiments confirmed the expected
	superlinear convergence behaviour for simple roots, while also illustrating the
	theoretical degradation to linear convergence in the presence of a multiple
	root.

\end{solution}

\begin{exercise}
	Write a function that performs inverse quadratic interpolation for finding a
	root of $f$, given three initial estimates $x_1,x_2,x_3$. Use this function
	to solve Exercise~\ref{ex:4.3.2} and determine the rate of convergence for
	the different roots.
\end{exercise}

\begin{solution}
	Inverse Quadratic Interpolation (IQI) is an iterative root-finding algorithm
	that fits a quadratic polynomial in $y$ to three prior data points
	$(x_{k-2}, y_{k-2})$, $(x_{k-1}, y_{k-1})$, and $(x_k, y_k)$, where $y =
		f(x)$. The root of this polynomial at $y=0$ provides the next iterate
	$x_{k+1}$.

	To solve the equations defined in Exercise~\ref{ex:4.3.2}, the
	following sets of initial estimates $(x_0, x_1, x_2)$ were strategically
	chosen from the domains:
	\begin{itemize}
		\item $f_1(x) = x^2 - e^{-x}$: initial estimates $(0.0, 0.5, 1.0)$
		\item $f_2(x) = 2x - \tan(x)$: initial estimates $(1.0, 1.2, 1.4)$
		\item $f_3(x) = e^{x+1} - 2 - x$: initial estimates $(-1.2, -0.5, 0.5)$
	\end{itemize}

	\begin{center}
		\centering
		\includegraphics[width=\textwidth]{figures/4-4/quadratic-interpolation-plot.pdf}
		\captionof{figure}{Plots of the functions $f_1(x)$, $f_2(x)$, and
			$f_3(x)$ showing the roots determined via inverse quadratic
			interpolation. The dashed line represents $y=0$, and the calculated
			roots are denoted by points.}
		\label{fig:iqi_plots}
	\end{center}

	\paragraph{Results}
	The algorithm successfully located the roots for all three functions, as
	visualised in Figure \ref{fig:iqi_plots}. The estimated orders of
	convergence ($q$) and asymptotic error constants ($C$) are presented in
	Table \ref{tab:iqi_results}.

	\begin{center}
		\centering
		\captionof{table}{Roots found via Inverse Quadratic Interpolation,
			alongside the estimated order of convergence ($q$) and asymptotic
			constant ($C$). Theoretical $q \approx 1.839$ for simple roots.}
		\label{tab:iqi_results}
		\renewcommand{\arraystretch}{1.2}
		\begin{tabular}{l c S[table-format=2.6] S[table-format=1.3] S[table-format=1.3]}
			\toprule
			Function & Initial Estimates & {Root ($x^*$)} & {$q$} & {$C$} \\
			\midrule
			$f_1(x)$          & $(0.0, 0.5, 1.0)$          & 0.703467                & 1.888 & 0.930 \\
			$f_2(x)$          & $(1.0, 1.2, 1.4)$          & 1.165561                & 1.842 & 3.890 \\
			$f_3(x)$          & $(-1.2, -0.5, 0.5)$        & -1.000000               & 1.020 & 0.672 \\
			\bottomrule
		\end{tabular}
	\end{center}

	\paragraph{Discussion on Rapid Convergence and Algorithm Selection}
	The theoretical order of convergence for IQI when applied to simple roots is
	$q \approx 1.839$.
	Our numerical results for $f_1(x)$ and $f_2(x)$ yield $q \approx 1.888$ and
	$q \approx 1.842$, respectively, which closely align with this theoretical
	expectation. For $f_3(x)$, the convergence rate degrades strictly to linear
	($q \approx 1.020$). As discussed in the previous exercise, $x = -1$ is a
	root of multiplicity $m=2$ for $f_3(x)$, which intrinsically forces
	higher-order methods (like Secant and IQI) to revert to linear convergence.

	To calculate the convergence rate $q$, only the least-squares fitting
	technique (Algorithm 1) was used. The successive differences
	approach (Algorithm 2) consistently failed for this exercise,
	producing \texttt{NaN} outputs and triggering warnings for insufficient
	iterations.

	This failure highlights a limitation of local error-estimation techniques:
	IQI converges so rapidly ($q \approx 1.839$ is faster than the Secant
	method's $1.618$) that the sequence of errors hits the machine epsilon
	in very few iterations. As a result, there are not enough valid, monotonically
	decreasing data points before catastrophic cancellation corrupts the sequence.
	The global, logarithmic regression of Algorithm 1 bypasses this local noise,
	making it the superior, and necessary, choice for evaluating rapidly converging
	algorithms.
\end{solution}

\subsection{Some observations about the algorithms used in this section}

The numerical experiments performed in this chapter highlight the different
trade-offs between robustness, convergence speed, and computational cost. A
useful quantity for understanding the sensitivity of a root-finding problem is
the condition number of the root:
\begin{equation*}
	\kappa_r = \left| f'(r) \right|^{-1}.
\end{equation*}
When $\left| f'(r) \right|$ is small, small perturbations in the function value
may correspond to large perturbations in the root location. This explains why
multiple roots, where the derivative vanishes at the root, are typically much
harder to approximate accurately.

The asymptotic error constant is defined as
\begin{equation*}
	C = \lim_{n \to \infty}
	\frac{\left| x_{n+1} - r \right|}
	     {\left| x_n - r \right|^q},
\end{equation*}
where $q$ is the order of convergence. In the asymptotic regime, a finite
positive value of $C$ indicates that the method follows the expected error law
$|x_{n+1}-r| \approx C |x_n-r|^q$.

To compare methods more fairly, one must also account for the work required by
each iteration. If an iterative method has convergence order $q$ and each
iteration requires $s$ units of work, its efficiency index is
\begin{equation*}
	E = q^{1/s}.
\end{equation*}
Here, a unit of work usually corresponds to one function evaluation, or to the
evaluation of a derivative such as $f'(x)$. The following summary reports the
metrics obtained from the numerical tests:

\begin{center}
	\centering
	\captionof{table}{Comparison of convergence rate, computational cost, and
		efficiency index for the root-finding methods considered in this chapter.}
	\label{tab:nonlinear_method_metrics}
	\renewcommand{\arraystretch}{1.2}
	\begin{tabular}{l c c c}
		\toprule
		Method & Convergence Rate & Computational Cost & Efficiency Index \\
		\midrule
		Newton          & 1.999767                  & 6                           & 1.414131                  \\
		Secant          & 1.607453                  & 8                           & 1.607453                  \\
		IQI             & 1.839000                  & 9                           & 1.839000                  \\
		Bisection       & 1.000000                  & 48                          & 1.000000                  \\
		\bottomrule
	\end{tabular}
\end{center}

\begin{center}
	\centering
	\captionof{table}{Additional convergence diagnostics for Newton's method.}
	\label{tab:nonlinear_newton_details}
	\renewcommand{\arraystretch}{1.2}
	\begin{tabular}{l c}
		\toprule
		Metric & Value \\
		\midrule
		Newton $C$ tail & $[0.333343,\ 0.352967,\ 0.353578]$ \\
		Synthetic $q$   & 2.000000                            \\
		Synthetic $C$   & 0.500000                            \\
		\bottomrule
	\end{tabular}
\end{center}

The data in Tables~\ref{tab:nonlinear_method_metrics} and
\ref{tab:nonlinear_newton_details} were obtained from the automated tests on the
reference problem $f(x)=x^2-2$, whose exact root is $\sqrt{2}$. The convergence
rates were estimated from the tail of the error sequence, so they reflect the
asymptotic behaviour of each method after the first transient iterations. The
reported computational cost corresponds to the number of iterates required in
the test, while the efficiency index accounts for the number of function or
derivative evaluations needed by one iteration. Newton's method therefore shows
nearly quadratic convergence, the secant method is close to the expected
superlinear order, IQI is assigned its theoretical order, and bisection remains
linear but requires the largest number of steps. The second table confirms that
the convergence estimator is reliable: on a synthetic sequence with prescribed
$q=2$ and $C=0.5$, it recovers both values accurately.

\paragraph{Bisection method}
The bisection method converges slowly, but its strength is that it requires only
continuity and a bracketing interval where the function changes sign. After
$n$ bisection steps, the root is confined to an interval $(a_n,b_n)$ of length
$2^{-n}(b_0-a_0)$. If the midpoint $m_n$ is used as an approximation of the
root $r$, then
\begin{equation*}
	|m_n-r| < 2^{-(n+1)}(b_0-a_0).
\end{equation*}
Thus, bisection is linearly convergent. It is reliable and independent of the
smoothness of $f(x)$, but it generally requires many iterations to reach high
accuracy.

\paragraph{Newton method}
For simple roots, Newton's method exhibits quadratic convergence
($q \approx 2$). This means that, once the iterates enter the asymptotic regime,
the number of correct digits roughly doubles at each step. Its speed comes at
the cost of requiring both $f(x)$ and $f'(x)$ at every iteration. Consequently,
Newton's method is very efficient for smooth, well-conditioned problems, but it
can lose its quadratic behaviour near multiple roots or when the derivative is
small.

When the multiplicity $m$ of a root is known, the modified Newton iteration
\begin{equation*}
	x_{k+1} = x_k - m \frac{f(x_k)}{f'(x_k)}
\end{equation*}
restores quadratic convergence under the usual regularity assumptions.

\paragraph{Interpolation methods}
The secant method replaces the derivative in Newton's method with a finite
difference approximation. It therefore avoids the explicit evaluation of
$f'(x)$, but its theoretical convergence rate for simple roots drops to
$q = (1+\sqrt{5})/2 \approx 1.618$, strictly between linear and quadratic
convergence.

Inverse Quadratic Interpolation extends this idea by fitting an inverse
quadratic model through three previous points. For simple roots it has a higher
theoretical convergence order, approximately $q \approx 1.839$, and in the tests
it reached the root very rapidly. As with the secant method, however, the
observed convergence can deteriorate near multiple roots because the underlying
root-finding problem becomes ill-conditioned.
